% ============================================================
%  鹰眼长空 — 程序设计说明书  排版导言区
%  编译：xelatex main.tex  (两遍，生成目录与交叉引用)
% ============================================================
\usepackage{geometry}
\geometry{left=2.6cm,right=2.6cm,top=2.6cm,bottom=2.4cm,headheight=15pt,footskip=1.0cm}

\usepackage{amsmath,amssymb,amsthm,bm}
\usepackage{graphicx}
\usepackage{float}
\usepackage{array,booktabs,longtable,tabularx,multirow,makecell}
\usepackage{xltabular}
\usepackage{enumitem}
\usepackage{xcolor}
\usepackage{tikz}
\usetikzlibrary{arrows.meta,positioning,shapes.geometric,shapes.misc,shapes.symbols,fit,backgrounds,calc,chains,decorations.pathreplacing,decorations.markings,patterns,matrix}
\usepackage{caption}
\usepackage{fancyhdr}
\usepackage{listings}
\usepackage[most]{tcolorbox}
\usepackage[hidelinks,bookmarksnumbered,bookmarksopen]{hyperref}
\usepackage{url}

% ---------- 黑白配色（白底黑字） ----------
\definecolor{bwtext}{gray}{0}
\definecolor{bwgray}{gray}{0.35}
\definecolor{bwlight}{gray}{0.92}
\definecolor{bwfill}{gray}{0.96}
\color{black}

% ---------- 字体（SimHei 无粗体，启用伪粗体） ----------
\setCJKmainfont[AutoFakeBold=2.4]{SimSun}
\setCJKsansfont[AutoFakeBold=2.4]{SimHei}
\setCJKmonofont[AutoFakeBold=2.0]{SimSun}

% ---------- 标题格式 ----------
\ctexset{
  section = {
    format = \Large\heiti,
    name = {,},
    number = \arabic{section},
    aftername = \hspace{0.7em},
    beforeskip = 1.2ex plus .2ex,
    afterskip = 1.0ex plus .2ex,
  },
  subsection = {
    format = \large\heiti,
    number = \arabic{section}.\arabic{subsection},
    aftername = \hspace{0.7em},
    beforeskip = 1.0ex plus .2ex,
    afterskip = 0.8ex plus .2ex,
  },
  subsubsection = {
    format = \normalsize\heiti,
    number = \arabic{section}.\arabic{subsection}.\arabic{subsubsection},
    aftername = \hspace{0.6em},
    beforeskip = 0.8ex plus .2ex,
    afterskip = 0.6ex plus .2ex,
  },
}

% ---------- 图表编号：章-序 ----------
\counterwithin{figure}{section}
\counterwithin{table}{section}
\counterwithin{equation}{section}
\renewcommand{\thefigure}{\thesection-\arabic{figure}}
\renewcommand{\thetable}{\thesection-\arabic{table}}
\renewcommand{\theequation}{\thesection-\arabic{equation}}
\captionsetup{font=small,labelsep=quad,justification=centering}
\captionsetup[table]{skip=4pt}
\captionsetup[figure]{skip=4pt}
\renewcommand{\figurename}{图}
\renewcommand{\tablename}{表}
\renewcommand{\lstlistingname}{代码}

% ---------- 页眉页脚 ----------
\pagestyle{fancy}
\fancyhf{}
\fancyhead[L]{\small\songti 鹰眼长空 — 低空无人机智能监测系统}
\fancyhead[R]{\small\songti 程序设计说明书}
\fancyfoot[C]{\small\thepage}
\renewcommand{\headrulewidth}{0.4pt}
\renewcommand{\footrulewidth}{0pt}
\fancypagestyle{plain}{\fancyhf{}\fancyfoot[C]{\small\thepage}\renewcommand{\headrulewidth}{0pt}}

% ---------- 列表 ----------
\setlist[itemize]{leftmargin=2em,itemsep=2pt,topsep=3pt,parsep=0pt}
\setlist[enumerate]{leftmargin=2.4em,itemsep=2pt,topsep=3pt,parsep=0pt}

% ---------- 代码清单 ----------
\lstdefinestyle{py}{
  language=Python,
  basicstyle=\ttfamily\footnotesize,
  keywordstyle=\bfseries,
  commentstyle=\color{bwgray},
  stringstyle=\color{black},
  frame=single,
  rulecolor=\color{black!40},
  framexleftmargin=10pt,
  numbers=left,
  numberstyle=\tiny\color{bwgray},
  numbersep=6pt,
  breaklines=true,
  breakatwhitespace=false,
  columns=fullflexible,
  keepspaces=true,
  showstringspaces=false,
  tabsize=4,
  aboveskip=6pt,
  belowskip=6pt,
  xleftmargin=22pt,
  extendedchars=true,
  inputencoding=utf8,
}
\lstdefinestyle{json}{
  basicstyle=\ttfamily\footnotesize,
  frame=single,
  rulecolor=\color{black!40},
  breaklines=true,
  columns=fullflexible,
  showstringspaces=false,
  aboveskip=6pt,
  belowskip=6pt,
}
\lstdefinestyle{sh}{
  language=bash,
  basicstyle=\ttfamily\footnotesize,
  frame=single,
  rulecolor=\color{black!40},
  breaklines=true,
  columns=fullflexible,
  showstringspaces=false,
  aboveskip=6pt,
  belowskip=6pt,
}
\lstset{style=py}

% ---------- 提示框（黑白） ----------
\newtcolorbox{notebox}[1][]{
  enhanced,
  colback=white,
  colframe=black!55,
  boxrule=0.5pt,
  arc=2pt,
  left=6pt,right=6pt,top=4pt,bottom=4pt,
  fonttitle=\bfseries\heiti,
  title=#1,
  breakable,
}

% ---------- 插图配色（正文仍为黑白） ----------
\definecolor{cA}{RGB}{222,235,249}   % 蓝
\definecolor{cB}{RGB}{228,243,231}   % 绿
\definecolor{cC}{RGB}{252,241,222}   % 橙
\definecolor{cD}{RGB}{238,240,245}   % 灰
\definecolor{cLine}{RGB}{58,78,104}  % 深蓝灰描边
\definecolor{cLbl}{RGB}{45,45,45}

% ---------- TikZ 统一风格 ----------
\tikzset{
  base/.style   = {draw=cLine, line width=0.5pt, fill=cA, align=center,
                   font=\small, inner sep=3pt, minimum height=8mm, text=black},
  proc/.style   = {base, rectangle, rounded corners=1.5pt, minimum width=26mm},
  deci/.style   = {base, diamond, aspect=2.2, inner sep=1pt, minimum width=20mm,
                   minimum height=9mm, font=\footnotesize, fill=cC},
  term/.style   = {base, rectangle, rounded corners=9pt, minimum width=24mm, fill=cB},
  data/.style   = {base, trapezium, trapezium left angle=72, trapezium right angle=108,
                   text width=26mm, inner sep=3pt, fill=cD},
  doc/.style    = {base, signal, signal to=east, minimum width=26mm, inner sep=5pt, fill=cD},
  layer/.style  = {draw=cLine, line width=0.7pt, fill=cD!35, inner sep=8pt,
                   rounded corners=3pt},
  grp/.style    = {draw=cLine!80, line width=0.5pt, rounded corners=2pt, inner sep=7pt, fill=white},
  arr/.style    = {-{Stealth[length=2.2mm]}, line width=0.5pt, draw=cLine},
  darr/.style   = {-{Stealth[length=2.2mm]}, line width=0.4pt, draw=cLine, dashed},
  busr/.style   = {-{Stealth[length=2.2mm]}, line width=1.0pt, draw=cLine},
  lbl/.style    = {font=\scriptsize, inner sep=1.5pt, fill=white, text=black},
}

% ---------- 自定义命令 ----------
\newcommand{\code}[1]{\texttt{\detokenize{#1}}}
\newcommand{\file}[1]{\path{#1}}   % 路径使用 url 宏包，可在 / 处断行
\def\UrlBreaks{\do\/\do\_\do\.\do\-}
\newcommand{\term}[1]{\textbf{#1}}
\renewcommand{\arraystretch}{1.25}
\setcounter{tocdepth}{2}
\renewcommand{\topfraction}{0.85}
\renewcommand{\bottomfraction}{0.7}
\renewcommand{\textfraction}{0.12}
\renewcommand{\floatpagefraction}{0.75}
\setlength{\parindent}{2em}
\setlength{\parskip}{0.35ex}
\setlength{\emergencystretch}{2.5em}
\tolerance=1500


% ---------- 截图标注点（第 11 章使用） ----------
\tikzset{
  hmark/.style={circle, draw=white, line width=0.9pt, fill=cLine,
               text=white, font=\bfseries\scriptsize, inner sep=1.4pt,
               minimum size=4.6mm, opacity=0.96, text opacity=1},
}
\newcommand{\hotspot}[3]{\node[hmark] at (#2,#3) {#1};}
\newcommand{\cn}[1]{\textcircled{\footnotesize #1}}
